\subsection{KT-TWOFIELD-GLOBAL: arbitrary Galois elements do not remove the example}
\label{sec:kt-twofield-global}

\begin{corollary}[KT-TWOFIELD-GLOBAL]
For the two forms in KT-TWOFIELD-661, let
$r_D:G_{\mathbf Q}\to\operatorname{GL}_2(
\overline{\mathbf Q}_{661})$ be the standard Galois
representation attached to $f_D$, with arithmetic-Frobenius
trace $a_q(f_D)$ and determinant
$\chi_D\epsilon^{536}$, where $\epsilon$ is the
$661$-adic cyclotomic character, and put
$\omega=\epsilon\bmod661$. Define
\[
 R_D=1\oplus(r_D\otimes\epsilon^{-268})|_{G_{E_D}},
 \qquad E_D=\mathbf Q(\sqrt{-D}).
                                                               \tag{TG.1}
\]
At the coefficient prime supplied by KT-TWOFIELD-661,
\[
 \overline R_D^{\mathrm{ss}}\simeq
          1\oplus\omega^{-268}\oplus\omega^{268}
                    \qquad(D=3,7).
                                                               \tag{TG.2}
\]
Hence the difference between $\operatorname{tr}R_D(\sigma)$
and $1+\epsilon(\sigma)^{-268}+\epsilon(\sigma)^{268}$
is nonunit for \emph{every} $\sigma\in G_{E_D}$.
Both $f_3$ and $f_7$ are non-CM.
\end{corollary}

\begin{proof}
The attached integral Galois representation and its trace
and determinant normalization are recalled in
Fretwell--Roberts, Section 2.1. The first paragraph of
the proof of their Theorem 3.5, equation (5), explicitly
applies Chebotarev and Brauer--Nesbitt to an Eisenstein
congruence and obtains its residual semisimplification.
Here it gives
\[
              \overline r_D^{\mathrm{ss}}
                     \simeq1\oplus\chi_D\omega^{536}.
                                                               \tag{TG.3}
\]
Equivalently, this follows by continuity from (TF.3) on
Frobenius classes and linear independence of simple
residual trace functions; $661>2$ prevents a multiplicity
ambiguity. Restrict (TG.3) to $G_{E_D}$, where $\chi_D=1$,
twist by $\omega^{-268}$, and add the trivial character.
This proves (TG.2) and the assertion for every $\sigma$.

At level three, KT-GL2-CM gives zero CM dimension because
$3\nmid268$. At level seven the same CM level formula
from Sch\"utt, Theorems 1.1 and 1.4, forces any CM newform
of this level to be induced from an unramified Hecke
character of $\mathbf Q(\sqrt{-7})$. The prime three is
inert in that field, so its Hecke eigenvalue would be
$a_3=0$. But (TF.3) gives
\[
                     a_3(f_7)\equiv1-3^{536}
                                      \equiv50\pmod{661},
                                                               \tag{TG.4}
\]
a contradiction. The last residue is a direct modular
exponentiation. Thus neither congruence can be removed
by deleting only CM newforms from the classical spaces.
\end{proof}

This assertion concerns the indicated classical residual
systems. It does not assert a new automorphic admissibility
or packet-coverage theorem. Whenever an established ambient
model includes these systems, its trace comparison has
this global obstruction. KT-TWOFIELD-661 still proves
$K_{536}(\mathbf Z)\{661\}=0$ independently: a genuine
global congruence is a necessary source condition in
the carried K17-S17-JOINT, not a sufficient condition for the
rational Selmer group to be nonzero.

